
\newpage
\TOCadd{Abstract}

\section*{\large Abstract}

Canada faces a growing number of abandoned oil wells and petrochemical contamination sites. The presence of petroleum hydrocarbons (PHC) in soil can pose serious risks to the environment, agricultural productivity, and human health in nearby communities if left unchecked. Remediating these sites is also costly and time-intensive, particularly for the many legacy and orphaned wells where responsible parties are difficult to identify or no longer solvent. In situ bioremediation technologies artificially accelerate biodegradation within the soil without requiring excavation of the affected material. Effective in situ biostimulation remediation depends on accurately mapping PHC plume distributions from sparse sensor measurements, as interpolating these measurements into a continuous plume estimate can significantly improve remediation efficiency and outcomes. This thesis evaluates two traditional distance-based spatial interpolation methods, Inverse Distance Weighting (IDW) and Kriging, alongside three parameter estimation techniques, 4D-Variational Data Assimilation (4D-Var), the Gauss-Levenberg-Marquardt algorithm (GLM), and the Iterative Ensemble Smoother (IES). The spatial interpolation methods are benchmarked on a physics-based synthetic dataset and validated on real-world sensor data to identify the most effective spatial interpolation method for petrochemical remediation in Canada. 

Across nearly all conditions tested, including measurement noise, sensor dropout, and reduced sensor counts, the physics-based parameter estimation methods (4D-Var, GLM, IES) substantially outperform the distance-based methods. Among these, IES achieves the best overall accuracy and additionally provides an ensemble-based uncertainty estimate, though at an increased computational cost. Kriging, meanwhile, offers no advantage over IDW, as its reliance on variogram estimation is poorly suited to the spatial sparsity of the problem. Taken together, these results indicate that where computational cost is acceptable, IES is the recommended method for PHC plume reconstruction at active remediation sites.